\section{Discussion}
\label{sec:discussion}

\paragraph{Two scaling axes saturate; the third does not.} \Cref{tab:scaling-curves} settles a common assumption that reasoning scaling and interaction scaling lie on a single shared compute curve --- ``you can always think more.'' They do not. Reasoning saturates at 73.3\% (86.7\% under the generous thinking-heavy partition), a ceiling fixed by the data-processing inequality: longer chains of thought reorganize information already present in the weights but add none. Best-of-$N$ sampling saturates at 86.7\%, drawing only on output-distribution variance and no new conditional-on-task information. Iteration with execution feedback reaches $\geq 97.8\%$ by injecting genuinely new information about the program's runtime behavior on each cycle. The first two axes are bounded; the third is not, and the boundary between them is the difference between recombining internal knowledge and importing external evidence.

\paragraph{Architecture buys efficiency and reliability.} Among the three interaction strategies --- single-agent loop (L), sample-and-select (IAD), and proposer--reviewer harness (H) --- the harness is the most token-efficient (1{,}029 vs.\ 1{,}431 L vs.\ 1{,}416 IAD mean output tokens) and the only variant that holds strict 100\% with zero seed variance. Three composing mechanisms produce this: \emph{context isolation} (the reviewer attends to the minimum information needed to localize defects rather than a multi-thousand-token generative history), \emph{structured critique} (a typed defect list rather than raw stderr), and \emph{role specialization} (in-context steering toward a diagnostic frame is sharper when uncontaminated by the generative history). The load-bearing comparison is \emph{interaction strategies vs.\ reasoning/sampling-only} --- a 13--27\,pp gap robust across seeds, model families, and held-out tasks --- and within the interaction tier the architectural separation of proposer and reviewer is what converts that gap into a deployable, variance-free system at the lowest token cost.

\paragraph{Grounding the evaluation, not just the feedback.} The visual modalities expose a principle that generalizes well beyond them: \emph{the quality of a rendered artifact is a geometric property, and a geometric property must be measured deterministically}. A vision--language model reading a screenshot is the wrong instrument. It cannot see text-on-text overlap, and it cannot see canvas overflow at all, because the overflowing content is cropped off-frame before the image reaches the judge --- the screenshot judge rates 14 of 15 dense academic-figure renders ``perfect'' when only 3 are clean. Measuring the rendered DOM directly --- every element's true bounding box, computing overlaps, clipping, and overflow exactly --- is exact where the screenshot is blind. Grounding \emph{both} the reviewer and the scorer in DOM geometry yields large, significant reductions in real layout defects across all four visual modalities measured identically --- academic figures $-74\%$, dense slides $-73\%$, responsive web pages $-47\%$, and SVG/CSS animations $-40\%$ (all two-sided sign-test $p<0.002$; \Cref{tab:geometric}, \Cref{sec:geometric}). Ungrounded \emph{evaluation} is as much a failure as ungrounded \emph{feedback}; grounding both is the principle, and the two-sided grounding is what recovers the full effect on visual artifacts.

\paragraph{Variance is a budget.} The RFT result (\Cref{sec:distill-recipes}) inverts the assumption that polishing the policy further is monotonically good. RFT improves every per-turn consistency metric we measure \emph{and} regresses pass@3 by 17\,pp. For deployment with sampling, output-distribution variance \emph{is} the inference-scaling lever: any post-SFT method that reduces variance --- expert iteration, distillation iterations, self-consistency training --- spends from the same budget that pass@$k$ draws on. Variance is not noise to be minimized; it is the resource that best-of-$N$ converts into quality.

\section{Conclusion}
\label{sec:conclusion}

Test-time compute has a third scaling axis. Interaction scaling is fundamentally distinct from reasoning and sampling because each interaction cycle injects information from outside the model's weights, and we formalize this through a Type 0/1/2/3 feedback taxonomy and a data-processing-inequality bound that pins reasoning- and sampling-only scaling below the interaction ceiling. The framework's three predictions --- a hierarchy of feedback channels ordered by actionability, reasoning-only saturation below the interaction ceiling, and architectural separation of proposer and reviewer buying efficiency and reliability --- are each tested and confirmed.

The results hold across modality, across model family, and on held-out tasks. A budget-aware proposer--reviewer harness over Claude Sonnet 4 lifts quality on every one of six modalities, with the execution-grounded code channel producing the largest and most reliable lift ($+33.3$\,pp to a strict $100\%$); an in-modality control that holds the task suite fixed and varies only the feedback type confirms it is executing the tests --- not the extra reviewing pass --- that carries the gain, exactly as the taxonomy predicts. At a 20K-token budget on hard code tasks, reasoning-only saturates at 73.3\% (86.7\% under the generous thinking-heavy partition) and best-of-$N$ at 86.7\%, while all three iteration-with-feedback strategies reach $\geq 97.8\%$; the harness is the most token-efficient (1{,}029 vs.\ 1{,}431 L vs.\ 1{,}416 IAD mean output tokens) and the only variant at strict 100\% across multiple seeds. The effect replicates across three model families (Sonnet $+33.3$\,pp, Qwen3-235B $+26.7$\,pp, GPT-5 $+13.3$\,pp on a higher single-shot baseline) and on a 32-task held-out code suite (3/3 recoveries, zero regressions), and the allocation simplex shows an 86.6\,pp spread, monotone in the proposer's per-call budget. For the visual modalities we ground the evaluation as well as the feedback: a deterministic DOM-geometry metric, applied to both reviewer and scorer, cuts real layout defects by $40$--$74\%$ across academic figures, dense slides, responsive web pages, and SVG/CSS animations (all $p<0.002$). And the harness's interaction-quality behavior distills into an 8B Qwen3-VL student via SFT on judge-filtered teacher trajectories: $0.50\times$ teacher mean@1 capability on a 44-task out-of-distribution suite (rising to $0.70\times$ at pass@2), 44\%/56\% pass@1/pass@3 on an 18-task hard held-out suite, at $\sim 10\times$ lower deployment cost.

The operational recipe is concrete. Wrap a frontier model in a proposer--reviewer harness with the most actionable grounded feedback channel available for the modality. For visual artifacts, ground \emph{both} the reviewer and the scorer in deterministic geometry, measuring the rendered DOM rather than a screenshot. Allocate at least half the per-task budget to the proposer, and spend extra inference compute on additional samples rather than deeper iteration. For cost-sensitive deployment, distill the proposer into a small markdown-emitting student and run it inside the same harness. Interaction scaling is real, distinct from reasoning and sampling, predictable from the feedback taxonomy, and immediately deployable.
